\documentclass{article}
\usepackage{preamble}

\begin{document}
\header{Filemon Mateus}{CS 6140}{Frequent Items}{Homework \#\ 5}{\today}
\begin{enumerate}[leftmargin={*}, font={\bf}, label={\arabic*.}, ref={\arabic*}]
  \item \label{qst:1}
    \begin{enumerate}[label={(\Alph*)}, ref={\theenumi(\Alph*)}]
      \item \label{qst:1A}
        Below is the output of running a fully fledge vanilla implementation of
        {\sf Misra Gries} on {\sf S1.txt}:
        \begin{table}[H]
          \centering
          \begin{tabular}{|l|c|c|c|c|c|c|c|c|c|c|c|}
            \hline
            {\sf\bfseries label} & {\sf ``a''} & {\sf ``b''} & {\sf ``p''} & {\sf ``d''} & {\sf ``s''} & {\sf ``f''} & {\sf ``c''} & {\sf ``v''} & {\sf ``k''} & {\sf ``t''} & {\sf ``n''} \\ \hline
            {\sf\bfseries count} & 772,865 & 472,863 & 2 & 1 & 1 & 1 & 233,850 & 2 & 1 & 1 & 1 \\ \hline
            {\sf\bfseries ratio} & 0.25762 & 0.15762 & 0.0 & 0.0 & 0.0 & 0.0 & 0.07795 & 0.0 & 0.0 & 0.0 & 0.0 \\ \hline
          \end{tabular}
          \caption{
            Output of {\sf Misra Gries} algorithm on stream {\sf S1}. See
            \autoref{lst:fast-misra-gries} for a fuller \\ explanation on
            how these frequencies have been generated.
          }
          \label{tbl:misra-gries-table-s1}
        \end{table}

        Below is the output of running a fully fledge vanilla implementation of
        {\sf Misra Gries} on {\sf S2.txt}:
        \begin{table}[H]
          \centering
          \begin{tabular}{|l|c|c|c|c|c|c|c|c|c|c|c|}
            \hline
            {\sf\bfseries label} & {\sf ``t''} & {\sf ``b''} & {\sf ``a''} & {\sf ``h''} & {\sf ``j''} & {\sf ``s''} & {\sf ``c''} & {\sf ``u''} & {\sf ``f''} & {\sf ``i''} & {\sf ``o''} \\ \hline
            {\sf\bfseries count} & 0 & 710,525 & 1,911,225 & 0 & 0 & 0 & 311,750 & 0 & 0 & 0 & 0 \\ \hline
            {\sf\bfseries ratio} & 0.0 & 0.17763 & 0.47781 & 0.0 & 0.0 & 0.0 & 0.07794 & 0.0 & 0.0 & 0.0 & 0.0 \\ \hline
          \end{tabular}
          \caption{
            Output of {\sf Misra Gries} algorithm on stream {\sf S2}. See
            \autoref{lst:fast-misra-gries} for fuller \\ explanation on how
            these frequencies have been generated.
          }
          \label{tbl:misra-gries-table-s2}
        \end{table}

        {\bf\sffamily --- Elements that ``might'' occur at least 25\% of the time} \par

        We note that for a fix entry, say $q$, in the stream, the following
        inequality holds between the ground-truth frequency $f_q$ and the
        estimated frequency $\hat{f}_q$:
        \[
          f_q - \frac{m}{k} \leq \hat{f}_q
        \]
        Choose $m=3,000,000$, $k=12$, and $f_q = 1/4 \times 3,000,000 = 750,000$
        to obtain that $\hat{f}_q \geq 500,000$. The task now reduces to finding
        all elements $q$ of the stream whose estimated frequency (corresponding
        to the count entries from \autoref{tbl:misra-gries-table-s1}) is greater
        or equal to $500,000$. The only element that satisfies this constraint
        on its estimated frequency is item: {\sf ``a''}. \\

        Similarly, for \autoref{tbl:misra-gries-table-s2}, choose $m=4,000,000$,
        $k=12$, and $f_q = 1/4 \times 4,000,000 = 1,000,000$ to obtain that
        $\hat{f}_q \geq 666,666.\bar{6}$. From \autoref{tbl:misra-gries-table-s2},
        the only elements satisfying this constraint are items {\sf ``a''} and
        {\sf ``b''}. \\

        {\bf\sffamily --- Elements that ``must'' occur at least 25\% of the time} \par
        
        Elements of {\sf S1} that are guaranteed to occur at least 25\% times must
        satisfy the following constraint:
        \[
          \hat{f}_q \geq f_q = 1/4 \times 3,000,000 = 750,000
        \]
        From \autoref{tbl:misra-gries-table-s1}, the only item that exceeds such
        frequency is item {\sf ``a''}. \\
        
        Similarly, the elements of {\sf S2} that must occur at least 25\% of the
        time must satisfy the following constraint:
        \[
          \hat{f}_q \geq f_q = 1/4 \times 4,000,000 = 1,000,000
        \]

        From \autoref{tbl:misra-gries-table-s2}, the only item that exceeds such
        frequency is item {\sf ``a''}.

      \item \label{qst:1B}
        Below is the output of running a fully fledge vanilla implementation of
        {\sf Count Min Sketch} on {\sf S1.txt}:
        \begin{table}[H]
          \centering
          \begin{tabular}{|l|c|c|c|}
            \hline
            {\sf\bfseries label} & {\sf ``a''} & {\sf ``b''} & {\sf ``c''} \\ \hline
            {\sf\bfseries count} & 899,566 & 659,805 & 360,551 \\ \hline
            {\sf\bfseries ratio} & 0.29986 & 0.21993 & 0.12018 \\ \hline
          \end{tabular}
          \caption{
            Output of {\sf Count Min Sketch} algorithm on stream {\sf S1}. See
            \autoref{lst:count-min-sketch} for a fuller \\ explanation on how
            these frequencies have been generated.
          }
          \label{tbl:count-min-sketch-s1}
        \end{table}

        Below is the output of running a fully fledge vanilla implementation of
        {\sf Count Min Sketch} on {\sf S2.txt}:
        \begin{table}[H]
          \centering
          \begin{tabular}{|l|c|c|c|}
            \hline
            {\sf\bfseries label} & {\sf ``a''} & {\sf ``b''} & {\sf ``c''} \\ \hline
            {\sf\bfseries count} & 2,000,100 & 839,401 & 400,625 \\ \hline
            {\sf\bfseries ratio} & 0.50003 & 0.20985 & 0.10016 \\ \hline
          \end{tabular}
          \caption{
            Output of {\sf Count Min Sketch} algorithm on stream {\sf S2}. See
            \autoref{lst:count-min-sketch} for a fuller \\ explanation on how
            these frequencies have been generated.
          }
          \label{tbl:count-min-sketch-s2}
        \end{table}

        {\bf\sffamily --- Elements that ``might'' occur at least 25\% of the time} \par
        
        From Section 11.2.3 in M4D, irrespective of the randomness incurred by the hash
        functions, the following invariance overestimate always holds:
        \[
          \hat{f}_q \geq f_q
        \]
        Therefore, for stream {\sf S1},  we want to enforce $\hat{f}_q \geq 1/4 \times
        3,000,000 = 750,000$ for items that ``might'' occur at least 25\% of the time.
        From \autoref{tbl:count-min-sketch-s1}, the only item satisfying this constraint
        is item {\sf ``a''}. \\

        Likewise, for stream {\sf S2} we want to enforce $\hat{f}_q \geq 1/4 \times
        4,000,000 = 1,000,000$ for items that ``might'' occur at least 25\% of the time.
        From \autoref{tbl:count-min-sketch-s2}, the only item satisfying this constraint
        is item {\sf ``a''}. \\

        {\bf\sffamily --- Elements that ``must'' occur at least 25\% of the time} \par
        
        Elements of {\sf S1} that are guaranteed to occur at least 25\% of the time
        must satisfy the following constraint:
        \[
          \hat{f}_q > f_q + \varepsilon F_1
        \]
        Choose $\varepsilon = 2/k = 1/6$ to obtain for stream {\sf S1} that $\hat{f}_q
        > 1,250,000$. Inspecting the values from \autoref{tbl:count-min-sketch-s1},
        item {\sf ``a''} is the only item satisfying this constraint. \\
        
        Analagously, for the same fixed choice of $\varepsilon = 1/6$, in stream {\sf
        S2} we obtain $\hat{f}_q > 1,666,666.6\bar{6}$. Item {\sf ``a''} is the only
        item satisfying this constraint with an estimated frequency of $2,000,100$.
        
      \item \label{qst:1C}
        {\sf Misra Gries} -- use a hash map with keys representing labels/words and
        values representing their corresponding indices in the counter. \\

        {\sf Count Min Sketch} -- compute each hashed word by summing up the {\sf ascii}
        values of its constituent characters and then use modular hashing. To circumvent
        the problem of having anagrams being hashed to the same location, multiply the
        {\sf ascii} values of the characters by the position their appear in the word.

      \item \label{qst:1D}
        While the space for the {\sf Count-Min Sketch} is larger than the {\sf Misra-Gries
        Sketch}, it has an advantage that its guarantees hold under the so-called turnstyle
        streaming model. Under this model, each element $a_i \in A$ can either add one or
        subtract one from corpus, but each count $f_j$ must remain non-negative. Thus {\sf 
        Count-Min} has the same guarantees in the turnstile model, but {\sf Misra-Gries}
        does not.
    \end{enumerate}
\end{enumerate}

\newpage

\section*{Appendix}

\lstinputlisting
  [caption={\sf fast-misra-gries.py}, label={lst:fast-misra-gries}]
  {dev-dir/week5/fast-misra-gries.py}
  
\newpage

\lstinputlisting
  [caption={\sf count-min-sketch.py}, label={lst:count-min-sketch}]
  {dev-dir/week5/count-min-sketch.py}
\end{document}
